\documentclass[11pt]{article}
\usepackage{fontspec,amsmath,amssymb,amsthm,hyperref}
\usepackage[margin=1in]{geometry}
\setlength{\emergencystretch}{3em}
\title{Finite quadratic transport transfers second moments}
\author{LangevinLean mathematical exposition}
\date{}
\begin{document}\maketitle
\noindent Source module: \texttt{UnadjustedLangevin/TransportMoments.lean}. Compile this proof document with XeLaTeX or LuaLaTeX. The accompanying HTML pairs the mathematical argument with the exact Lean statements and records the scope of the formalization.


Finite \(W_2\) distance to a finite-second-moment law supplies the other second moment automatically. This removes a redundant integrability premise when passing an initial law to the L\ensuremath{{}^{2}} formulation.


\section{memLp id of wasserstein2 ne top}


Finite Wasserstein distance supplies a finite-cost coupling. On that coupling space the second coordinate is in \(L^2\) by the reference moment, and the coordinate difference is in \(L^2\) by finite coupling cost. Their sum is the first coordinate, so its law \(\mu\) has a finite second moment. No optimal coupling is selected.



\noindent Lean declarations: \texttt{memLp\_id\_of\_wasserstein2\_ne\_top}.

\section{wasserstein2 ne top of memLp}


Use the explicit independent coupling \(\mu\otimes\pi\). Each coordinate is in \(L^2\), hence their difference is in \(L^2\) and has finite cost. Wasserstein distance is at most this coupling cost, proving finiteness.



\noindent Lean declarations: \texttt{wasserstein2\_ne\_top\_of\_memLp}.

\section{memLp of law of wasserstein2 ne top}


The finite-transport theorem first proves \(\mathrm{id}\in L^2(\mu)\). Replace \(\mu\) by the supplied pushforward law of \(X\), then use the change-of-law \(L^2\) criterion to pull this moment back to \(X\) on its original sample space.



\noindent Lean declarations: \texttt{memLp\_of\_law\_of\_wasserstein2\_ne\_top}.

\section{memLp of wasserstein2 map ne top}


Specialize the preceding theorem to \(\mu=P\circ X^{-1}\), whose law identity is reflexive. Finite transport of the actual pushforward to a finite-moment reference therefore gives \(X\in L^2(P)\).



\noindent Lean declarations: \texttt{memLp\_of\_wasserstein2\_map\_ne\_top}.

\end{document}
